\documentclass[12pt,a4paper,oneside]{book}

\usepackage[scheme=chinese,fontset=none]{ctex}

\setCJKmainfont{LXGWWenKai}[
  Path           = ./fonts/,
  Extension      = .ttf,
  UprightFont    = LXGWWenKai-Regular,
  BoldFont       = LXGWWenKai-Medium,
  ItalicFont     = LXGWWenKai-Light,
  BoldItalicFont = LXGWWenKai-Medium,
  Scale          = 1.0
]
\setCJKsansfont{LXGWWenKai}[
  Path         = ./fonts/,
  Extension    = .ttf,
  UprightFont  = LXGWWenKai-Regular,
  BoldFont     = LXGWWenKai-Medium,
  Scale        = 1.0
]
\setCJKmonofont{LXGWWenKaiMono}[
  Path        = ./fonts/,
  Extension   = .ttf,
  UprightFont = LXGWWenKaiMono-Regular,
  BoldFont       = LXGWWenKaiMono-Medium,
  ItalicFont     = LXGWWenKaiMono-Light,
  BoldItalicFont = LXGWWenKaiMono-Medium,
  Scale       = 0.92
]

\setmainfont{LXGWWenKai}[
  Path        = ./fonts/,
  Extension   = .ttf,
  UprightFont = LXGWWenKai-Regular,
  BoldFont    = LXGWWenKai-Medium,
  ItalicFont  = LXGWWenKai-Light,
  BoldItalicFont = LXGWWenKai-Medium,
  Scale       = 1.0
]
\setsansfont{LXGWWenKai}[
  Path        = ./fonts/,
  Extension   = .ttf,
  UprightFont = LXGWWenKai-Regular,
  BoldFont    = LXGWWenKai-Medium,
  Scale       = 1.0
]
\setmonofont{LXGWWenKaiMono}[
  Path        = ./fonts/,
  Extension   = .ttf,
  UprightFont = LXGWWenKaiMono-Regular,
  BoldFont       = LXGWWenKaiMono-Medium,
  ItalicFont     = LXGWWenKaiMono-Light,
  BoldItalicFont = LXGWWenKaiMono-Medium,
  Scale       = 0.92
]

\usepackage[margin=2.5cm,headheight=15pt]{geometry}

\XeTeXlinebreakskip = 0pt plus 0.15\ccwd minus 0.06\ccwd

\usepackage{mathtools}
\usepackage{amsmath,amssymb}
\usepackage{amsthm}
\usepackage{mathrsfs}
\usepackage{bbm}
\usepackage{bm}
\usepackage{esint}
\usepackage{siunitx}
\theoremstyle{plain}
\newtheorem{definition}{定义}[chapter]
\newtheorem{theorem}{定理}[chapter]
\newtheorem{proposition}{命题}[chapter]
\newtheorem{corollary}{推论}[chapter]
\newtheorem{lemma}{引理}[chapter]
\newtheorem{axiom}{公理}[chapter]
\newtheorem{property}{性质}[chapter]

\theoremstyle{definition} 
\newtheorem{example}{例}[chapter]

\theoremstyle{remark}  
\newtheorem{remark}{注}[chapter]

\usepackage{tabularx}
\usepackage{multirow}
\usepackage{booktabs}
\usepackage{array}
\usepackage{longtable}

\usepackage{makecell}
\usepackage[table]{xcolor}

\usepackage{graphicx}
\usepackage{adjustbox}

\newenvironment{fittikz}
  {\begin{adjustbox}{max width=\textwidth,center}}
  {\end{adjustbox}}
\usepackage{tikz}
\usepackage{tikz-cd}
\usepackage{pgfplots}
\usepackage{tikz-3dplot}
\usepackage{placeins}
\pgfplotsset{compat=1.18}
\usetikzlibrary{
  arrows,
  arrows.meta,
  shapes,
  shapes.geometric,
  decorations.pathmorphing,
  positioning,
  calc,
  patterns,
  decorations.pathmorphing,
  decorations.markings,
  decorations.pathreplacing,
  backgrounds,
  fit,
  matrix,
  scopes,
  shadows,
  through,
  quotes,
  angles
}

\usepackage{listings}
\lstset{
  basicstyle=\ttfamily\small,
  breaklines=true,
  frame=single,
  numbers=left,
  xleftmargin=2em,
  xrightmargin=2em
}

\usepackage{float}

\usepackage{afterpage}

\usepackage{titlesec}
\usepackage{titletoc}
\usepackage{fancyhdr}

\usepackage{placeins}
\usepackage{etoolbox}

\pretocmd{\section}{\FloatBarrier}{}{}
\pretocmd{\subsection}{\FloatBarrier}{}{}

\pagestyle{fancy}
\fancyhf{}
\fancyhead[L]{\leftmark}
\fancyhead[R]{}
\fancyfoot[C]{\thepage}
\renewcommand{\headrulewidth}{0.4pt}
\renewcommand{\footrulewidth}{0pt}

\titleformat{\chapter}[display]
  {\normalfont\huge\bfseries\centering}
  {第\,\thechapter\,章}
  {20pt}
  {\Huge}
\titleformat{name=\chapter,numberless}[display]
  {\normalfont\huge\bfseries\centering}
  {}
  {20pt}
  {\Huge}
\titlespacing*{\chapter}{0pt}{-30pt}{40pt}

\contentsmargin{2.55em}
\titlecontents{chapter}
  [0pt]
  {\addvspace{10pt}\bfseries}
  {第\thecontentslabel 章\quad}
  {}
  {\titlerule*[0.5pc]{.}\contentspage}

\titlecontents{section}
  [0pt]
  {\addvspace{2pt}}
  {\thecontentslabel\quad}
  {}
  {\titlerule*[0.5pc]{.}\contentspage}

\titlecontents{subsection}
  [0pt]
  {\addvspace{1pt}}
  {\thecontentslabel\quad}
  {}
  {\titlerule*[0.5pc]{.}\contentspage}

\titlecontents{figure}
  [0pt]
  {\addvspace{2pt}}
  {\thecontentslabel\quad}
  {}
  {\titlerule*[0.5pc]{.}\contentspage}

\titlecontents{table}
  [0pt]
  {\addvspace{2pt}}
  {\thecontentslabel\quad}
  {}
  {\titlerule*[0.5pc]{.}\contentspage}

\usepackage{enumitem}
\usepackage{caption}
\usepackage{subcaption}

\usepackage[protrusion=false]{microtype}
\setlength{\emergencystretch}{2em}

\usepackage[colorlinks=true,linkcolor=blue,citecolor=blue,urlcolor=cyan,unicode=true]{hyperref}

\hypersetup{
    pdftitle        = {周期数学},
    pdfauthor       = {Eloy Kim (eloy.kim@outlook.com)},
    pdfsubject      = {周期数学——以周转圆级数为统一语言，建立跨天文学、物理学、生物学、语言学、经济学与计算机科学的周期结构诊断框架。},
    pdfkeywords     = {周期数学; 周转圆; 调和分析; 守恒律; 对称性; 相位; 振幅; 频率; 周期卷积; 群表示论; 傅里叶变换; 梅林变换; 小波变换; 易经; 莱布尼茨; 历法; 芯片; 信号处理;},
    pdfcreator      = {Eloy Kim},
    pdfproducer     = {XeLaTeX with hyperref},
    pdfdisplaydoctitle = true,
    pdflang         = {zh-CN},
    unicode         = true,
    colorlinks      = true,
    linkcolor       = blue,
    citecolor       = blue,
    urlcolor        = cyan
}

\usepackage[capitalise,nameinlink]{cleveref}
\crefname{property}{性质}{性质}
\Crefname{property}{性质}{性质}
\crefname{axiom}{公理}{公理}
\Crefname{axiom}{公理}{公理}

